% ************************** Thesis Abstract *****************************
\begin{abstract}
This dissertation measures how floating-point precision, hardware, compiler
semantics and small implementation choices affect a stand-alone CPU/CUDA
high-resolution shock-capturing (HRSC) finite-volume code, and what those
effects mean for reproducing a published algorithm.

Part~I treats the compressible Euler equations with a nominally second-order
MUSCL--Hancock scheme and the HLLC Riemann solver, exercised on five one- and
two-dimensional cases: Sod, Toro3 and Toro5 against exact Riemann references,
and Liska--Wendroff configurations~3 and 12 against \(1600^2\) and \(800^2\)
fp64 numerical references. Matched strict-IEEE HLLC CPU and GPU saved outputs
agree exactly, at \(L_1=0\), \(L_\infty=0\) and \(\mathrm{ULP}_{\max}=0\) over
final states and checkpoints in all five cases. Direct fp32--fp64 differences
stay below the available reference or discretisation scale, the largest
density reference-scaled ratio reaching \(1.30\times10^{-4}\). Sensitivity is
axis-dependent: \texttt{-O2} and \texttt{-O3} are bit-identical where tested,
whereas fast-math perturbs non-stationary saved states.

Part~II carries the same questions into ideal magnetohydrodynamics (MHD) with
generalised Lagrange multiplier divergence cleaning and the
Harten--Lax--van~Leer (HLL) and HLL with Discontinuities (HLLD) solvers. A
circularly polarised Alfv\'en wave, the one case here with an exact solution,
verifies an observed order of \(1.818\) on both devices and both precisions.
Every reported run keeps density and pressure finite and positive. Across
Orszag--Tang and Kelvin--Helmholtz refinements, all 12 mean-\(L_1\) density
discrepancies between 32-bit (fp32) and 64-bit (fp64) calculations remain below
3\% of the matched fp64 finest adjacent-grid difference; in the domain mean
that difference is two to seven binary32 roundings on the coarse grids but
several hundred by \(512^2\), so refinement widens the precision gap rather
than closing it. Saved HLL CPU and GPU states are bit-identical, but only
because device multiply--add contraction was disabled to match the host
default: restoring the compiler's own default removes the agreement in every
pair and, on Brio--Wu in fp32, moves the stored density further than halving
the working precision does. That agreement is a property of the declared build
recipe, not of the hardware, whereas changing the baseline instruction set
changes nothing. Compiler floating-point semantics and the HLLD branch rule
both produce non-zero discrepancies, and the branch result has a structural
cause: HLL's wave-speed clamp leaves its equality branches unreachable, while
HLLD's stay live at a decision that instrumentation finds exactly tied at
\(62\%\) of Brio--Wu interfaces. Followed in time, the precision and build
axes all grow as power laws with exponents near 1.4--1.6, not exponentially.
Robustness over the tested horizons is therefore distinct from arithmetic
insensitivity, and reproducibility requires complete records of the numerical
path and execution environment.
\end{abstract}
